% Separate post-CP34 candidate. No promotion or least-index conclusion.
\subsection{Six both-large rows in the tracefree first layer}

Throughout this subsection, $K$ is complete discretely valued of mixed
characteristic $(0,5)$, its residue field is algebraically closed,
$e_K=20$, and $\mu _5\subset K$. We consider the remaining actual
$125$-torsion comparison of two whole potentially supersingular
height-two groups, with different finite special characters and
nontrivial residual wild action. Use the nonsquare finite-determinant
normalization of B77 and B91. Put $L/K$ for the common tame good field
of degree24, so $v_L(5)=480$. Valuations on every point field are
normalized to have value group $\mathbf Z$. A conductor below means
last upper break, without adding one. We restrict to the both-large
case
\[
 B_i=5b_i>20,\qquad b_i\in\{5,7,9,11,13,15,17,19\},
 \qquad b_1<b_2,
 \qquad (\dim W_1,\dim W_2)=(3,4).                 \tag{B-RD94.1}
\]
The actual high quotients $Q_i=M_i/H$ have a common $25$-comparison
and parameters $5j_i$, where $j_i=24-b_i$. Their old first-level
images are their common high line. No $125$-comparison of the $Q_i$
is assumed.

\begin{lemma}[The actual tracefree Iwahori group and its two point fields]
\label{lem:B-RD94-group}
\status{Candidate; complete independent audits pending.}
Let $F/L$ be the common original $125$-division field. In an integral
tame eigenbasis with original high line $\mathbf F_5e_1$, its group is
\[
 W=\left\{\begin{pmatrix}1+5a&v\\5w&1+5d\end{pmatrix}:
 \begin{array}{l}a,d,w\in\mathbf Z/25,\ v\in\mathbf Z/125,\\
 a_0+d_0=v_0w_0\text{ in }\mathbf F_5\end{array}\right\},
 \qquad |W|=5^8.                                    \tag{B-RD94.2}
\]
Here $a=a_0+5a_1$, $d=d_0+5d_1$, $w=w_0+5w_1$, and
$v=v_0+5v_1+25v_2$, with digits in $\{0,1,2,3,4\}$.
Let $F_2$ be the original $25$-division field, $F_Q$ the actual
$25$-division field of the high quotients, and $E=F_2F_Q$.
Then
\[
 [F_2:L]=[F_Q:L]=5^4,\quad [E:L]=5^5,\qquad
 \operatorname{Gal}(F/E)=
 \left\{1+25\begin{pmatrix}a_1&v_2\\0&d_1\end{pmatrix}\right\}.
                                                        \tag{B-RD94.3}
\]
The original and quotient residual fields $R,R'$ are the independent
cyclic degree-five fields of $v_0,w_0$. Moreover
$K=K_0(\mu _{25})$ for a Cohen coefficient field $K_0$, and the
actual cyclotomic degree-five extension $L(\mu _{125})/L$ has break576.
\end{lemma}
\begin{proof}
B91 gives $W_1=\mathfrak {sl}_2(\mathbf F_5)$ and $W_2=M_2(\mathbf F_5)$.
Thus the determinant modulo25 vanishes on the first congruence
subgroup. Its remaining homomorphism factors through the original
residual unipotent group. That group has a nontrivial tame weight,
whereas the determinant is invariant under conjugation. Hence this
homomorphism also vanishes. This puts $W$ in (B-RD94.2). The containing
group has exactly $5^8$ elements, the order already supplied by B91,
so equality holds. This is an equality of actual images.

The determinant on $L$ is the actual cyclotomic character, or its
inverse in the dual convention; its finite tame factor is trivial
there. Consequently $\mu _{25}\subset L$. A tame extension cannot
supply a missing five-primary cyclotomic extension, so
$\mu _{25}\subset K$. The Cohen-field and index argument proved in
B86 gives $K=K_0(\mu _{25})$. The next cyclotomic extension has
lower and upper break24 over $K$, and tame base change of degree24
gives576 over $L$. These statements use the full arbitrary-perfect-
residue argument already recorded in B86, not a finite-residue-field
substitution.

Modulo25 the original matrix is specified by
$(v_0,w_0,a_0,v_1)$, with $d_0=v_0w_0-a_0$. In the actual quotient
lattice $(e_1/5,e_2)$ its matrix is
\[
       \begin{pmatrix}1+5a&5v\\w&1+5d\end{pmatrix}\pmod {25},
                                                        \tag{B-RD94.4}
\]
which is specified by $(v_0,w_0,a_0,w_1)$. This proves all field
degrees and the exact kernel in (B-RD94.3). Both $v_0$ and $w_0$
are additive on the full group, and their joint image is
$\mathbf F_5^2$, proving the assertion about $R,R'$.
\end{proof}

\begin{lemma}[Normality orders the actual diagonal conductors]
\label{lem:B-RD94-ordering}
\status{Candidate; complete independent audits pending.}
Write $r=b(R/L)$ and $r'=b(R'/L)$. Over $R$, the three coordinate
characters of $\operatorname{Gal}(F_2/R)=\mathfrak {sl}_2(\mathbf F_5)$
are $w_0,a_0,v_1$, with conductors $\rho,x,C$, respectively.
Over $R'$, the opposite three characters of
$\operatorname{Gal}(F_Q/R')$ are $v_0,d_0,w_1$, with conductors
$\sigma,y,C'$. They satisfy
\[
 \rho<x<C,\qquad \sigma<y<C',\qquad
 x>r,\quad y>r',\quad 24\mid x,y,\quad
 y=x+4(r-r').                                      \tag{B-RD94.5}
\]
The actual degree-25 point field over either residual field has
its last ramification line equal to its canonical high-translation
line. The residual breaks satisfy
\[
 r\equiv-6b_i\pmod {24},\quad
 r'\equiv6b_i\pmod {24},\qquad
 r<600-30b_i,\quad r'<600-30j_i.                    \tag{B-RD94.6}
\]
\end{lemma}
\begin{proof}
For $H=\operatorname{Gal}(F_2/R)$, conjugation by a residual upper
unipotent of parameter $t$ acts on a tracefree matrix by
\[
 (a_0,v_1,w_0)\longmapsto
 (a_0+tw_0,\ v_1-2ta_0-t^2w_0,\ w_0).             \tag{B-RD94.7}
\]
Conductors are invariant under this conjugation. Taking $t=1,-1$
and adding and subtracting the corresponding characters gives
$\rho\le x\le C$. The tame weights of the three coordinates are
pairwise distinct: the diagonal weight is one and the two root
weights are opposite primitive fourth roots in $\mathbf F_5$.
Thus the inequalities are strict. Here is the precise reason that
no equal-conductor case is lost. If two independent cyclic
characters of different tame weights had the same conductor, their
actual elementary-abelian compositum would have one ramification
break. The injective leading-coefficient map of that ramification
group identifies tame conjugation with multiplication by one scalar
in the residue field. It cannot have two different eigenvalues in
$\mathbf F_5$ on two nonzero eigenvectors. The same argument in the
opposite quotient group proves $\sigma<y<C'$.

The point field $L(P_{25})$ for $5P_{25}=e_2$ has degree125:
the second-column orbit modulo25 is specified freely by
$(v_0,v_1,d_0)$. It contains $R$, and its extension over $R$ is
Galois of degree25, with translation characters $(v_1,d_0)$.
On $H$, $d_0=-a_0$, so their conductors are $C>x$; its last
line is therefore the $v_1$-line, exactly the original high line.
The same proof for (B-RD94.4), with the basis reordered, applies
to the quotient point field over $R'$. These are actual point
fields of the available $25$-modules.

We give a subgroup argument for $x>r$ to avoid assigning breaks
to arbitrary lifts. Let $D$ be the lower break in $H$ of any
matrix with $w_0=0$, $a_0\ne0$. Such a matrix may have any upper
coordinate; since $\rho<x<C$, its break is precisely $D$ and
$\phi_H(D)=x$. In a nontrivial coset of $H$ choose an element
$g$ with maximal lower break $M$. The quotient-average formula
gives $r=\phi_H(M)$: the other coset breaks are the minima of
$M$ and the corresponding subgroup breaks, including the term
$M$ itself. If $r\ge x$, then $M\ge D$. Choose $h\in H$ with
$w_0(h)\ne0$. Formula (B-RD94.7) says that $[g,h]$ has zero
lower coordinate and nonzero diagonal coordinate; hence its lower
break is $D$. But
\[
                    b([g,h])\ge b(g)+b(h)>D,          \tag{B-RD94.8}
\]
a contradiction. For completeness, this commutator inequality
follows directly from B60: both terms of
$(g-1)(h-1)\varpi-(h-1)(g-1)\varpi$ have valuation at least
$1+b(g)+b(h)$, and composing with $g^{-1},h^{-1}$ preserves
valuation. Thus no new ramification theorem is required. The
opposite argument gives $y>r'$.

Put $S=RR'$. Its actual cyclic restrictions have conductors
$\rho,\sigma$, and elementary Herbrand conversion gives
$\rho-\sigma=5(r'-r)$. Because $x>\rho$ and $y>\sigma$, the
two diagonal restrictions to $S$ have conductors
$5x-4\rho$ and $5y-4\sigma$. On the subgroup over $S$,
$v_0=w_0=0$ and the determinant relation reads $a_0=-d_0$.
They are therefore the same nonzero character up to sign. Equating
these two actual restricted conductors gives the last equality in
(B-RD94.5).

Choose a tame generator acting on a uniformizer of $L$ by $\zeta$.
The normalized high eigenvalue is $\zeta^{c_i-5b_i}$ and the
finite determinant is $\zeta^{6c_i}$, with $c_i\equiv1\pmod4$ and
$c_i\equiv5b_i\pmod6$. Thus its eigenvalue ratio is
$\zeta^{-4c_i-10b_i}=\zeta^{-6b_i}$. The quotient has the
opposite root weight. A tame eigen-uniformizer in a degree-five
residual field has eigenvalue $\zeta^5$; the injective leading
map gives exactly the two congruences in (B-RD94.6). A diagonal
cyclic character has weight one, giving $24\mid x,y$ in the same
way. Finally an original low first-level coordinate has normalized
valuation $5b_i$ over $R$, and each nontrivial translation has
displacement $600-25b_i$. The strict divisible-valuation part of
B60 gives $r<600-30b_i$. The quotient calculation gives the other
strict inequality. This proves every assertion with the actual
base field indicated.
\end{proof}

\begin{lemma}[The needed common-high point-field estimates]
\label{lem:B-RD94-common-high}
\status{Candidate; complete independent audits pending.}
In either of the two actual degree-25 point fields over $R$ or
$R'$ above, suppose its two parameter values are $m_1<m_2$.
Its normalized coordinates and translation displacements are
\[
 A_i=5m_i,\qquad D_{0,i}=125m_i,\qquad
 D_{1,i}=15000-625m_i,\qquad v(5)=60000.             \tag{B-RD94.9}
\]
Let $t<z$ be its lower breaks; its last line is the common high
line by the preceding lemma.
If both $m_i$ are prime to five, then $t=24m_1$.
For any of
\[
                 (m_1,m_2)=(5,9),(5,13),(5,17),(15,19),
                                                        \tag{B-RD94.10}
\]
one instead has
\[
 0<t<24m_1,\qquad24\mid t,\qquad
 z=15000-630m_2-4t,\qquad
 \operatorname{cond}(\text{off-diagonal})=3000-126m_2.
                                                        \tag{B-RD94.11}
\]
\end{lemma}
\begin{proof}
The original low inverse valuations are $m_i/25$ over $L$, and
the translating values are $m_i$ and $120-5m_i$. The actual point
field has degree125 over $L$, proving (B-RD94.9). The same
calculation applies to the actual quotient $25$-module. All changes
being used are strictly below the mixed-characteristic cutoff60000.
The full fixed-coefficient expansion of B60/B87 therefore applies.
For a coordinate with $5\nmid m$, the first possible change at
break $h$ is
$\min(5m+5h,J+h)$, where $J>5m$ is its first nonzero prime-to-five
exponent, possibly infinite. With $r_m=24m$, the two strict cases
force
\[
\begin{array}{c|c|c}
 &J&\text{relation at the last break }z\\\hline
 t<r_m&5m+4t&z+4t=15000-630m\\
 t>r_m&125m-t&z-t=15000-750m.
\end{array}                                           \tag{B-RD94.12}
\]
In the first row cancellation below $D_0$ forces the displayed
coefficient; at $z>t$ its ordinary term is uniquely first. The
second row is likewise an exact first-term equality.

If both parameters are prime, the two same strict rows in
(B-RD94.12) force equality of the parameters. The mixed strict
case would give $t=150m_1-126m_2<24m_1$, opposite to its assumed
position. At $t=24m_2$, model1 forces
$z=15000-750m_1+24m_2$; model2 requires
$J_2\ge101m_2$, whereas its last change would force
$J_2=750m_1-649m_2$. This would give $m_1\ge m_2$.
Only $t=24m_1$ remains. This repeats the numerical proof of
Lemma B92's common-high assertion, whose proof uses exactly the
data (B-RD94.9), without importing its surrounding crossing scope.

Now take one of (B-RD94.10). The second parameter is prime to five,
while $5\mid m_1$. If $t>24m_2$, every divisible-exponent term
of the first coordinate changes strictly after $D_{0,1}$:
its leading exponent $A_1$ is divisible by25, and every exponent
of five-adic order one is strictly larger than $A_1$. Thus its
first ordinary exponent is $125m_1-t$. The second row for model2
gives $z=t+15000-750m_2$, and, for all four listed pairs,
$5z+A_1>D_{1,1}$. All divisible terms are therefore later than
the required last displacement; the ordinary term forces
$z-t=15000-750m_1$, a contradiction. If $t=24m_2$, the first
coordinate still has that same ordinary exponent and forces
$z=15000-750m_1+24m_2$. The second coordinate's boundary argument
in the preceding paragraph again gives $m_1\ge m_2$.

Hence $t<24m_2$, and model2 gives the displayed value of $z$.
For these four pairs the two elementary inequalities
\[
 60000+630m_1-3630m_2>0,\qquad 630m_2-750m_1>0          \tag{B-RD94.13}
\]
hold. The first implies
$5z+A_1>D_{1,1}$, so model1 has a unique first ordinary exponent
$J_1=D_{1,1}-z$. At the first break its change is later than
$D_{0,1}$, since
\[
             J_1+t-D_{0,1}=630m_2-750m_1+5t>0.
\]
If $t\ge24m_1$, all divisible-exponent terms would also change
strictly after $D_{0,1}$: their first possible degree-five change
is at least $A_1+5t$, with equality impossible at the degree
$A_1$ because it is divisible by25. This is a contradiction.
Thus $0<t<24m_1$. Its cyclic quotient killing the high line is the
actual diagonal character, so $24\mid t$. The other upper break
is $t+(z-t)/5=3000-126m_2$, as asserted.
\end{proof}

\begin{proposition}[Four exclusions and eighteen remaining field profiles]
\label{prop:B-RD94-profiles}
\status{Candidate; complete independent audits pending.}
Under (B-RD94.1), the parameter pairs
\[
             (b_1,b_2)=(5,13),(5,17),(7,19),(11,19)   \tag{B-RD94.14}
\]
are impossible. For either $(5,9)$ or $(15,19)$, the actual field
$E=F_2F_Q$ has one of the following nine profiles. Write $(R_0,X)$
for one of
\[
 (18,24),(18,48),(18,72),(18,96),
 (42,48),(42,72),(42,96),(66,72),(66,96).              \tag{B-RD94.15}
\]
Its five upper breaks are
\[
 6,quad R_0,quad\frac{X+4R_0}{5},quad126,quad
                   \frac{1866+4R_0}{5}.              \tag{B-RD94.16}
\]
Its five lower breaks, denoted $\ell_1,\ldots,\ell_5$, are
\[
\begin{split}
 \ell_1&=6,\qquad \ell_2=5R_0-24,\qquad\ell_3=5X-24,\\
 \ell_4&=15726-20X-100R_0,\\
 \ell_5&=170226-20X+400R_0.
\end{split}                                           \tag{B-RD94.17}
\]
In $(5,9)$ one has $(r,r',x)=(R_0,6,X)$; in $(15,19)$ one has
$(r,r',y)=(6,R_0,X)$. The actual restricted cyclotomic character
on $\operatorname{Gal}(F/E)$ has conductor
\[
                  C_{\mathrm{cyc}}=803976-20X-2100R_0.
                                                        \tag{B-RD94.18}
\]
\end{proposition}
\begin{proof}
For $(5,13)$ or $(5,17)$, the mixed lemma gives $0<x<120$ and
$24\mid x$. The quotient parameter pairs are $(11,19)$ and
$(7,19)$, whose prime-parameter result gives respectively
$y=264,168$. In both cases $r'<30$ and $r'\equiv6\pmod{24}$,
so $r'=6$, while $r\equiv18\pmod{24}$. Equation (B-RD94.5)
gives $x=288-4r$ or $x=192-4r$. The only possible positive
multiple of24 below120 is respectively $(r,x)=(66,24)$ or
$(42,24)$. But then
$\sigma=5r-24=306$ or186, contradicting $\sigma<y$.
Interchanging the two actual $25$-division groups gives exactly
the other two pairs of (B-RD94.14). This interchange uses their
$25$-modules and point fields only, never quotient $125$-torsion.

For $(5,9)$, the quotient pair is $(15,19)$. Both instances of
the mixed lemma apply. Thus $0<x<120$, $0<y<360$,
$24\mid x,y$. The quotient residual break is $r'=6$, while
$r\equiv18\pmod{24}$ and $y=x+4r-24$. Together with $x>r$
these give exactly (B-RD94.15), with $(R_0,X)=(r,x)$.
The last off-diagonal conductors over the respective residual
fields are $C=1866$ and $C'=606$. Over $R$ the three conductors
are $6,X,1866$; over $R'$ they are
$5R_0-24,X+4R_0-24,606$.
The coset-maximum quotient formula, or its equivalent subgroup
Herbrand conversion, now gives for $F_2/L$ the upper breaks
\[
 6,R_0,(X+4R_0)/5,(1866+4R_0)/5,
\]
and for $F_Q/L$ the upper breaks
$6,R_0,(X+4R_0)/5,126$.

We justify precisely the passage to their compositum. Its two
kernel lines are the independent coordinates $v_1,w_1$.
Above126, upper quotient compatibility forces the group to be
the $v_1$-line, until the larger displayed last break. Just below126
it contains that line and projects nontrivially to the $w_1$-line,
so has both lines. At each preceding break its projections to the
common degree-$5^3$ quotient and these already-present kernel
lines determine its order. Thus its orders drop successively by
five at exactly the union of the two lists; there is no possible
additional drop with unchanged projections. This proves
(B-RD94.16). Integrating slopes $1,5,25,125,625$ gives (B-RD94.17).

For $(15,19)$, reverse the original and quotient computations:
now the original residual break is6, the quotient parameter pair
is $(5,9)$, and the smaller diagonal conductor is $y=X$.
The same nine profiles and the same compositum breaks result.
In this case the original $v_1$-character occurs at the fourth
break, whereas $w_1$ occurs at the fifth; their positions are
reversed in $(5,9)$.

The last upper break in (B-RD94.16) is at most426, strictly below
576. Therefore the actual cyclotomic subgroup is still nontrivial
after the $E$-quotient has become trivial. Upper quotient
compatibility and subgroup conversion give its exact conductor
on $\operatorname{Gal}(F/E)$ as
\[
 \psi_{E/L}(576)
 =\ell_5+3125\left(576-\frac{1866+4R_0}{5}\right),
\]
which expands to (B-RD94.18). This is a statement about the
actual cyclotomic subfield of $F$, not a separately chosen cyclic
extension.
\end{proof}

\begin{lemma}[Bounds from the two actual primitive $125$-point fields]
\label{lem:B-RD94-terminal}
\status{Candidate; complete independent audits pending.}
On the actual terminal group in (B-RD94.3), the characters
$a_1,d_1$ satisfy the following bounds for each profile
$(R_0,X)$ in (B-RD94.15):
\[
\begin{array}{c|r|r}
 (b_1,b_2)&\operatorname{cond}_E(a_1)\le&
                  \operatorname{cond}_E(d_1)\le\\\hline
 (5,9)&489096+80X-1600R_0&326976+80X+400R_0\\
 (15,19)&327096+80X+400R_0&489576+80X-1700R_0.
\end{array}                                           \tag{B-RD94.19}
\]
\end{lemma}
\begin{proof}
Take $P_h^+,P_\ell^+$ to be the two basis vectors of the actual
$125$-module in the chosen integral tame eigenbasis. Put
$P_h=5P_h^+$, $P_\ell=5P_\ell^+$, so that
$5P_h=e_1$ and $5P_\ell=e_2$ at the first level.
Write
\[
 E_h=L(P_h),\quad E_\ell=L(P_\ell),\quad
 E_h^+=L(P_h^+),\quad E_\ell^+=L(P_\ell^+),\quad
 S_h=RE_h,\quad T_h=RE_h^+.
\]
The full actual matrix group gives
\[
\begin{gathered}
 [E_h:L]=25,\ [E_\ell:L]=125,\ [S_h:L]=125,\\
 [E_h^+:L]=625,\quad[T_h:L]=[E_\ell^+:L]=3125,\\
 \operatorname{Gal}(T_h/S_h)=\mathbf F_5^2(a_1,w_1),\qquad
 \operatorname{Gal}(E_\ell^+/E_\ell)=\mathbf F_5^2(v_2,d_1).
\end{gathered}                                        \tag{B-RD94.20}
\]
Here and below the parentheses specify the actual matrix-coordinate
homomorphisms. To verify normality rather than merely the degrees,
the subgroup for $S_h$ has $v_0=a_0=w_0=0$; its quotient killing
the stabilizer of $P_h^+$ has additive coordinates $(a_1,w_1)$.
The subgroup for $E_\ell$ has $v_0=v_1=d_0=0$, hence $a_0=0$;
its quotient killing the stabilizer of $P_\ell^+$ has additive
coordinates $(v_2,d_1)$. Products give these assertions directly
modulo125.

Moreover $E/S_h$ is elementary abelian of degree25 with coordinates
$(v_1,w_1)$ and lower breaks $\ell_4,\ell_5$. Its upper breaks
are $\ell_4,(\ell_5+4\ell_4)/5$. The extension $E/E_\ell$ is
cyclic of degree25: its matrices have lower-root coordinate
$w_0+5w_1$, whose fifth power is the terminal $w_1$ direction.
Its lower breaks are $(6,\ell_4)$ in the first row of
(B-RD94.19), and $(\ell_2,\ell_5)$ in the second row.
These are subgroup lower breaks read from (B-RD94.17).
The actual intersections are
$T_h\cap E=S_h(w_1)$ and $E_\ell^+\cap E=E_\ell$;
the coordinate stabilizers give both, although only the indicated
restriction maps are needed for the ensuing inequalities.

For high targets, the actual inverse Newton polygons are single
segments because $B_i>20$. Their constant valuations are $5j_i$
and then $j_i/5$, while the degree-five coefficient has value
$20b_i$; this vertex is strictly above the relevant degree-25
chord. The complete coefficient estimate already proved in B49
and B69 puts every other vertex above it as well. Thus
\[
             v_L(P_h)=j_i/5,\qquad v_L(P_h^+)=j_i/125.
                                                        \tag{B-RD94.21}
\]
In $T_h$, the latter coordinate has valuation $25j_i$.
A nonzero element of its high translation line $w_1=0$ displaces
it by $15625j_i$. B60 therefore gives the unconditional bound
\[
                b(h)\le15600j_i\quad(w_1(h)=0,h\ne1).
                                                        \tag{B-RD94.22}
\]
This bound remains valid although the coordinate valuation is
five-divisible.

Put $z_h=\operatorname{cond}_{S_h}(w_1)$. If the high line is
last, the first lower break is $z_h$, so the $a_1$ conductor is
at most $z_h+(15600j_{\min}-z_h)/5$. If that line is not last,
the $w_1$ quotient already has the largest break and the $a_1$
conductor is at most $z_h$. The same latter conclusion holds
for one break. In both rows $15600j_{\min}>z_h$, so uniformly
\[
 \operatorname{cond}_{S_h}(a_1)
       \le(15600j_{\min}+4z_h)/5.                    \tag{B-RD94.23}
\]
In the first row $j_{\min}=15$ and $z_h=\ell_4$; in the second
$j_{\min}=5$ and $z_h=(\ell_5+4\ell_4)/5$. Restricting these
actual characters to $E$ and using the stated upper breaks of
$E/S_h$ gives respectively
\[
 \operatorname{cond}_E(a_1)\le1170000-4\ell_5,
 \qquad
 \operatorname{cond}_E(a_1)\le390000-4\ell_4.        \tag{B-RD94.24}
\]
The arguments of both Herbrand functions lie above their last
break; this follows immediately from the nine pairs (B-RD94.15).

For low targets, the two successive inverse coordinates have
$v_L(P_\ell)=b_i/25$ and $v_L(P_\ell^+)=b_i/625$.
The latter has valuation $5b_i$ in $E_\ell^+$. Every group element
with $d_1\ne0$ translates outside the canonical high line, hence
has displacement $3125b_i$. Thus every element of a nontrivial
$d_1$-coset has lower break at most $3120b_i$. The exact quotient
average proves
\[
               \operatorname{cond}_{E_\ell}(d_1)
                         \le3120b_{\min}.            \tag{B-RD94.25}
\]
Restrict to $E$ through the actual cyclic degree-25 extensions
above. Their upper pairs are
$(6,(\ell_4+24)/5)$ and
$(\ell_2,(\ell_5+4\ell_2)/5)$.
Using $b_{\min}=5$ and15 respectively gives
\[
 \operatorname{cond}_E(d_1)\le389880-4\ell_4,
 \qquad
 \operatorname{cond}_E(d_1)\le1170000-4\ell_5-20\ell_2.
                                                        \tag{B-RD94.26}
\]
Again the arguments exceed the last upper breaks. Substitution
of (B-RD94.17) gives the table. All bounds concern the same
characters $a_1,d_1$ on the actual group $\operatorname{Gal}(F/E)$.
No disjointness or unrelated cyclic-quotient identification was
assumed, and only the original actual125 module was used to form
the last two point fields.
\end{proof}

\begin{theorem}[Six tracefree-first-layer rows are excluded]
\label{thm:B-RD94-six-rows}
\status{Candidate; complete independent audits pending.}
Under (B-RD94.1), none of
\[
 (b_1,b_2)=(5,9),(5,13),(5,17),(7,19),(11,19),(15,19)
                                                        \tag{B-RD94.27}
\]
can occur. Consequently the remaining pairs for layer shape$(3,4)$
are among
\[
                       (7,15),(9,13),(9,17),(11,15).  \tag{B-RD94.28}
\]
\end{theorem}
\begin{proof}
Four rows have already been excluded. In each of the other two,
the actual determinant on $\operatorname{Gal}(F/E)$ is
\[
                        1+25(a_1+d_1)\pmod{125}.
\]
Its conductor is at most the larger of the two bounds in
(B-RD94.19). But its exact cyclotomic conductor is (B-RD94.18).
For $(5,9)$ the two positive differences are
\[
 314880-100X-500R_0,\qquad477000-100X-2500R_0;
\]
for $(15,19)$ they are
\[
 476880-100X-2500R_0,\qquad314400-100X-400R_0.
\]
Every expression is strictly positive for every pair
(B-RD94.15), a contradiction concerning the same actual character.

For clarity, the necessary different-character pairs prior to this
exclusion are exactly
\[
\begin{gathered}
 (5,9),(5,13),(5,17),(7,15),(7,19),\\
 (9,13),(9,17),(11,15),(11,19),(15,19).
\end{gathered}
\]
Indeed $b_1\equiv b_2\pmod4$, and the conditions
$c_i\equiv1\pmod4$, $c_i\equiv5b_i\pmod6$ and
$c_1-5b_1=c_2-5b_2\pmod{24}$ list these pairs after removing
equal unordered special orbits $\{c_i,5c_i\}$. The only other
same-congruence pairs $(7,11),(13,17)$ give equal special orbits;
they were never unresolved different-character cases. Removing
(B-RD94.27) leaves (B-RD94.28).
\end{proof}

\begin{remark}[Exact scope and retained finite checks]
\label{rem:B-RD94-scope}
\status{Candidate; complete independent audits pending.}
This is a conditional exclusion inside layer shape$(3,4)$ of the
both-large residual-wild case. The four pairs (B-RD94.28), every
layer$(4,4)$ pair, and the irreducible-residual branch remain open.
No remaining row is asserted to be realizable. The unrestricted
least index at $p=5,e=20$ remains in $\{2,3\}$.

The invoked B60 fixed-coefficient proof, B49 complete coefficient
argument, B69 complete actual quotient/Newton proof, B86 complete
Cohen/cyclotomic proof, B91 layer reduction, B92 numerical
common-high proof, and B93 actual-field conductor argument were
reread in their retained canonical forms.
No new original-source opening is claimed in this candidate.
The standalone exact program
\path{code/reduction_descent/verify_both_large_tracefree_rows.py}
checks the entire finite matrix group and the displayed field,
coefficient and conductor arithmetic. Its finite scope does not
replace the normality, field, or ramification proofs. The exact
executed code, complete output, initial failed census version and
its correction, and the earlier explicitly unproved exploratory
note remain retained. The checker enumerates all 390,625 matrices,
12 actual point/division-field image maps, 125 terminal determinants,
31,250 relative-quotient products, 25 powers of the actual cyclic
lower-root subgroup, and 625 tracefree conjugations. It independently
reconstructs the four mixed coefficient rows and all 18 conductor
profiles. The corrected author run passed at 04:25:12.551765 UTC
on 20 November 2026 in 1.2287049839942483 seconds. The initial run
failed only in its final redundant character census, whose loop
started the second odd parameter at an even integer; every preceding
matrix and conductor check had passed. Its exact source, failure
output and stdout are preserved, and the loop correction is retained.
No initial failed output was overwritten without a preserved copy.
\end{remark}
